% xelatex add_tour_analysis_report.tex
\documentclass[11pt, a4paper, oneside]{article}

\usepackage{fontspec}
\usepackage{polyglossia}
\setdefaultlanguage{russian}
\setotherlanguage{english}
\setmainfont{Liberation Serif}
\setsansfont{Liberation Sans}
\setmonofont[Scale=0.82]{DejaVu Sans Mono}

\usepackage[
  a4paper,
  top=22mm, bottom=22mm,
  left=28mm, right=22mm,
  headheight=14pt
]{geometry}

\usepackage[protrusion=false,expansion=false]{microtype}
\usepackage{parskip}
\usepackage{booktabs}
\usepackage{array}
\usepackage{tabularx}
\usepackage{longtable}
\usepackage{enumitem}
\usepackage{graphicx}
\usepackage{float}
\usepackage{caption}
\usepackage{amsmath}

\usepackage[dvipsnames,svgnames,table]{xcolor}
\definecolor{AccentBlue}  {RGB}{25,  90, 160}
\definecolor{AccentRed}   {RGB}{165,  20,  20}
\definecolor{AccentOrange}{RGB}{190,  70,   0}
\definecolor{AccentGreen} {RGB}{28,  130,  55}
\definecolor{AccentGray}  {RGB}{85,   90,  100}
\definecolor{LightGray}   {RGB}{246, 248, 252}
\definecolor{RowRed}      {RGB}{255, 235, 235}
\definecolor{RowOrange}   {RGB}{255, 245, 225}
\definecolor{RowBlue}     {RGB}{235, 242, 255}
\definecolor{MidGray}     {RGB}{120, 125, 138}
\definecolor{DarkGray}    {RGB}{35,   38,  48}
\definecolor{CodeBg}      {RGB}{245, 246, 252}

\usepackage[
  unicode, colorlinks=true,
  linkcolor=AccentBlue,
  urlcolor=AccentBlue,
  pdftitle={Анализ процесса add\_tour — отчёт о рисках},
]{hyperref}

\usepackage{tcolorbox}
\tcbuselibrary{breakable,skins}

\newtcolorbox{critbox}[1][]{
  colback=RowRed, colframe=AccentRed,
  colbacktitle=AccentRed, coltitle=white,
  fonttitle=\sffamily\small\bfseries,
  boxrule=0.5pt, arc=3pt, breakable,
  top=5pt, bottom=5pt, left=8pt, right=8pt, #1}

\newtcolorbox{warnbox}[1][]{
  colback=RowOrange, colframe=AccentOrange,
  colbacktitle=AccentOrange, coltitle=white,
  fonttitle=\sffamily\small\bfseries,
  boxrule=0.5pt, arc=3pt, breakable,
  top=5pt, bottom=5pt, left=8pt, right=8pt, #1}

\newtcolorbox{infobox}[1][]{
  colback=LightGray, colframe=AccentBlue!50,
  colbacktitle=AccentBlue, coltitle=white,
  fonttitle=\sffamily\small\bfseries,
  boxrule=0.5pt, arc=3pt, breakable,
  top=5pt, bottom=5pt, left=8pt, right=8pt, #1}

\newtcolorbox{summarybox}{
  colback=RowBlue, colframe=AccentBlue,
  boxrule=0.6pt, arc=4pt,
  top=6pt, bottom=6pt, left=10pt, right=10pt}

\usepackage{fancyhdr}
\pagestyle{fancy}
\fancyhf{}
\fancyhead[L]{\small\color{MidGray}\sffamily add\_tour — анализ рисков}
\fancyhead[R]{\small\color{MidGray}\sffamily 14 марта 2026}
\fancyfoot[C]{\small\color{MidGray}\thepage}
\renewcommand{\headrulewidth}{0.4pt}
\renewcommand{\headrule}{\hbox to\headwidth{\color{AccentBlue}\leaders\hrule height 0.4pt\hfill}}
\renewcommand{\footrulewidth}{0pt}

\usepackage{titlesec}
\titleformat{\section}
  {\large\bfseries\sffamily\color{AccentBlue}}
  {\thesection.}{0.7em}{}
  [\vspace{1pt}{\color{AccentBlue}\rule{\linewidth}{0.5pt}}\vspace{3pt}]
\titleformat{\subsection}
  {\normalsize\bfseries\sffamily\color{DarkGray}}
  {\thesubsection.}{0.6em}{}
\titleformat{\subsubsection}
  {\small\bfseries\color{AccentOrange}}
  {}{0pt}{}

\titlespacing*{\section}{0pt}{14pt}{6pt}
\titlespacing*{\subsection}{0pt}{10pt}{3pt}
\titlespacing*{\subsubsection}{0pt}{7pt}{2pt}

\newcolumntype{L}{>{\raggedright\arraybackslash}X}
\renewcommand{\arraystretch}{1.35}
\setlength{\parskip}{5pt}
\setlength{\parindent}{0pt}
\setlist{nosep, leftmargin=1.4em, topsep=3pt}
\captionsetup{font=small, labelfont=bf}

\newcommand{\coderef}[1]{{\footnotesize\color{MidGray}\ttfamily #1}}
\newcommand{\critical}{\textcolor{AccentRed}{\textbf{критично}}}
\newcommand{\warning}{\textcolor{AccentOrange}{\textbf{предупреждение}}}
\newcommand{\info}{\textcolor{AccentBlue}{замечание}}

% ════════════════════════════════════════════════════════════════
\begin{document}

{\noindent\large\bfseries\sffamily\color{AccentBlue} add\_tour — анализ рисков}
\hfill{\small\color{MidGray} 14 марта 2026 · 🔴~4 · ⚠️~5 · ℹ️~3}
\vspace{1mm}
{\color{AccentBlue}\rule{\textwidth}{0.5pt}}
\vspace{4mm}

% ════════════════════════════════════════════════════════════════
\section{Критичные проблемы}

% ── Проблема 1 ──────────────────────────────────────────────────
\begin{critbox}[title={🔴 \#1 — Версионирование контейнеров не откатывается при сбое алгоритма}]

\textbf{Фаза:} Фаза 2 → Фаза 3 (переход к алгоритму маршрутизации)\\
\coderef{run\_add\_tour\_task() → move\_visited\_containers\_to\_new\_tour() → шаг D4}

\textbf{Описание.}
При переносе контейнеров в новый маршрут выполняется группа операций:
\begin{enumerate}
  \item Старые версии контейнеров закрываются (заполняются поля \texttt{next\_version\_date})
  \item Создаются новые версии с привязкой к новой дате/маршруту
  \item Переносятся зоны транспортировщиков
  \item Создаются планы для перенесённых контейнеров
  \item Новые версии привязываются к планам
\end{enumerate}
Все эти операции выполняются \emph{до} отправки запроса к алгоритму маршрутизации.
Если алгоритм возвращает ошибку, \texttt{execution\_log.compensate()}
откатывает шаги только частично: удаляет созданные объекты,
но \textbf{не восстанавливает поля версионирования}
(\texttt{previous\_version\_date}, \texttt{next\_version\_date})
у записей, изменённых на шаге~1.

\textbf{Последствие.} В базе остаются «закрытые» старые версии контейнеров
(поле \texttt{next\_version\_date} заполнено), тогда как соответствующие
новые версии удалены. Контейнеры теряют непрерывность версионной цепочки.

\textbf{Рекомендация.}
Шаг \texttt{move\_visited\_containers\_to\_new\_tour} уже логируется
с \texttt{action="complex"} и сохраняет \texttt{value\_before}.
Необходимо убедиться, что в \texttt{value\_before} включены
\textbf{все изменённые поля} (\texttt{next\_version\_date},
\texttt{previous\_version\_date}, \texttt{plan\_id})
старых версий, а не только ID.
\end{critbox}

% ── Проблема 2 ──────────────────────────────────────────────────
\begin{critbox}[title={🔴 \#2 — Статусы исходных планов посещений не откатываются}]

\textbf{Фаза:} Фаза 2 → Фаза 3\\
\coderef{run\_add\_tour\_task() → update\_plans\_status\_moved + update\_empty\_plans\_status\_moved}

\textbf{Описание.}
После переноса контейнеров планы посещений, из которых ушли контейнеры,
переводятся в статус «Перенесено». Это два отдельных шага журнала:
\texttt{update\_plans\_status\_moved} и \texttt{update\_empty\_plans\_status\_moved}.

Компенсация через \texttt{execution\_log.compensate()} должна откатывать
эти статусы при сбое алгоритма. Однако \texttt{compensate()} использует
\texttt{qs.update(**state)}, где \texttt{state} содержит \textbf{одно значение status\_id}
для всей группы планов. Если до переноса планы имели \emph{разные} статусы,
откат восстановит \textbf{неправильный} статус части планов.

\textbf{Последствие.} Планы посещений отображаются пользователю
в статусе «Перенесено», хотя фактического переноса не произошло.

\textbf{Рекомендация.}
При логировании \texttt{update\_plans\_status\_moved} сохранять
индивидуальный статус каждого плана до изменения
(формат \texttt{value\_before} уже поддерживает это через \texttt{prepare\_rows}).
\end{critbox}

% ── Проблема 3 ──────────────────────────────────────────────────
\begin{critbox}[title={🔴 \#3 — Журнал выполнения зависает при ранних ошибках API}]

\textbf{Фаза:} Фаза 0 (API-слой)\\
\coderef{WasteSiteVisitedContainers.add\_tour() → AddTourTaskExecutionLog.objects.create()}

\textbf{Описание.}
Журнал выполнения (\texttt{AddTourTaskExecutionLog}) создаётся
\emph{сразу после} получения или создания маршрута — до проверок бизнес-правил
(зоны РО, совместимость с ТС, подрядчик).
При провале любой из этих проверок управление передаётся в блок компенсации
(\texttt{ERR0}), который возвращает ошибку API, но \textbf{не вызывает
\texttt{execution\_log.mark\_failed()}}.

Журнал остаётся в статусе \texttt{running} бессрочно.

\textbf{Последствие.} Накопление «зависших» записей в таблице
\texttt{statistic\_add\_tour\_task\_execution\_log}
со статусом \texttt{running}. При мониторинге или повторном запуске
они могут ошибочно расцениваться как активные операции.

\textbf{Рекомендация.}
В блоке обработки ошибок \texttt{ERR0} добавить:
\texttt{execution\_log.mark\_failed("Ошибка проверки: ...")}
перед возвратом ответа API.
\end{critbox}

% ── Проблема 4 ──────────────────────────────────────────────────
\begin{critbox}[title={🔴 \#4 — Сбой в самой компенсации не обрабатывается}]

\textbf{Фаза:} Компенсация при сбое алгоритма\\
\coderef{add\_tour\_task() → execution\_log.compensate() → ERR2}

\textbf{Описание.}
При сбое алгоритма вызывается \texttt{execution\_log.compensate()},
которая должна откатить все шаги и передать управление в \texttt{ERR2}
для финального логирования.

Вызов \texttt{compensate()} \textbf{не обёрнут в \texttt{try-except}}.
Если \texttt{compensate()} сама завершится с исключением
(например, при обращении к недоступной таблице или неизвестной модели),
исключение поднимается выше и блок \texttt{ERR2} не выполняется.

\textbf{Последствие.}
Данные остаются в промежуточном состоянии без записи в журнал.
Ни уведомления пользователю, ни логирования ошибки не происходит.

\textbf{Рекомендация.}
Обернуть вызов \texttt{compensate()} в \texttt{try-except}:
\begin{quote}\ttfamily\small
try:\\
\quad execution\_log.compensate()\\
except Exception as e:\\
\quad logger.exception("Ошибка компенсации: \%s", e)\\
\quad execution\_log.mark\_failed(str(e))\\
raise
\end{quote}
\end{critbox}

% ════════════════════════════════════════════════════════════════
\section{Предупреждения}

% ── Предупреждение 1 ────────────────────────────────────────────
\begin{warnbox}[title={⚠️ \#5 — Обновление зоны РО пропускается при наличии контейнеров}]

\textbf{Фаза:} Фаза 1\\
\coderef{add\_tour\_task() → ветка B1="Нет" → update\_tour\_ro\_zone}

\textbf{Описание.}
Обновление зоны РО (\texttt{ro\_zone}) у маршрута выполняется
\emph{только в ветке}, когда нет контейнеров для добавления без плана.
Если контейнеры есть — ветка с обновлением \texttt{ro\_zone} полностью пропускается.

\textbf{Вопрос к команде:} Должна ли \texttt{ro\_zone} обновляться независимо
от наличия контейнеров? Если да — это баг.
\end{warnbox}

% ── Предупреждение 2 ────────────────────────────────────────────
\begin{warnbox}[title={⚠️ \#6 — Polygon-only сценарий не закрывает журнал выполнения}]

\textbf{Фаза:} Фаза 0 (специальный сценарий)\\
\coderef{add\_tour() → polygon-only ветка → OK0}

\textbf{Описание.}
В «только полигонном» сценарии (обработка точек без переноса контейнеров)
API возвращает успех синхронно.
Если на этом пути успел создаться \texttt{AddTourTaskExecutionLog},
он не закрывается через \texttt{mark\_completed()}.

\textbf{Рекомендация.}
Убедиться что в polygon-only ветке либо журнал не создаётся вовсе,
либо явно закрывается перед возвратом ответа.
\end{warnbox}

% ── Предупреждение 3 ────────────────────────────────────────────
\begin{warnbox}[title={⚠️ \#7 — Отсутствует обработка ошибок при подготовке данных для алгоритма}]

\textbf{Фаза:} Фаза 3\\
\coderef{BuildRoutesPlanCommand.handle() → шаг D2 (подготовка: задания, точки, ТС)}

\textbf{Описание.}
Узел подготовки данных для алгоритма (задания, транспортные базы,
водители, точки) не имеет явной ветки ошибки в графе процесса.
При сбое на этапе подготовки поведение зависит от реализации
\texttt{BuildRoutesPlanCommand.handle()} — не ясно, будет ли вызван
\texttt{execution\_log.compensate()}.

\textbf{Рекомендация.}
Добавить явный \texttt{try-except} вокруг подготовки данных
с переходом к компенсации при ошибке.
\end{warnbox}

% ── Предупреждение 4 ────────────────────────────────────────────
\begin{warnbox}[title={⚠️ \#8 — RoutePlanningCalculation остаётся без статуса ошибки}]

\textbf{Фаза:} Фаза 3\\
\coderef{run\_add\_tour\_task() → RoutePlanningCalculation.objects.create() → сбой алгоритма}

\textbf{Описание.}
При сбое алгоритма запись \texttt{RoutePlanningCalculation}
остаётся в базе со статусом, не отражающим ошибку.
Поле \texttt{error\_message} модели существует, но не заполняется.

\textbf{Рекомендация.}
В блоке компенсации обновлять:
\texttt{RoutePlanningCalculation.objects.filter(pk=calc.id)\\
.update(status="failed", error\_message=str(e))}.
\end{warnbox}

% ── Предупреждение 5 ────────────────────────────────────────────
\begin{warnbox}[title={⚠️ \#9 — M2M-поля зон транспортировщиков не откатываются стандартным механизмом}]

\textbf{Фаза:} Фаза 4\\
\coderef{create\_and\_link\_waste\_site\_visit\_plans() → plan.transporter\_zones.set(...)}

\textbf{Описание.}
Планы посещений получают зоны транспортировщиков через
\texttt{plan.transporter\_zones.set(...)}.
Метод \texttt{execution\_log.compensate()} использует
\texttt{QuerySet.update()} для отката, который \textbf{не работает с M2M-полями}.

\textbf{Вопрос к команде:} Предусмотрен ли отдельный механизм отката
для \texttt{transporter\_zones}?
\end{warnbox}

% ════════════════════════════════════════════════════════════════
\section{Замечания}

% ── Замечание 1 ─────────────────────────────────────────────────
\begin{infobox}[title={ℹ️ \#10 — Дублирование проверки наличия переноса}]

Проверка «были контейнеры на перенос?» выполняется дважды:
в начале Фазы~2 (C3) и в начале Фазы~5 (F1).
Второй раз это не ошибка, но создаёт точку рассогласования при рефакторинге.
Можно передавать флаг \texttt{had\_containers\_to\_move}
из Фазы~2 в Фазу~5 явно.
\end{infobox}

% ── Замечание 2 ─────────────────────────────────────────────────
\begin{infobox}[title={ℹ️ \#11 — Два терминала успеха в одном процессе}]

Процесс имеет два независимых успешных завершения:
«только полигоны» (синхронный ответ API) и основной поток
(асинхронная задача, \texttt{mark\_completed()}).
Это архитектурно обоснованно, но стоит документировать явно —
чтобы мониторинг и тесты учитывали оба пути.
\end{infobox}

% ── Замечание 3 ─────────────────────────────────────────────────
\begin{infobox}[title={ℹ️ \#12 — Точка входа не связана с графом переходов}]

В схеме процесса узел «E1 Пользователь запускает add\_tour» изолирован:
нет ребра к следующему шагу.
Это артефакт инструмента построения схемы, не баг в коде.
Стоит исправить схему для корректной документации.
\end{infobox}



\vfill
\hrule
\vspace{4mm}
{\small\color{MidGray}
  Отчёт сгенерирован инструментом \texttt{поток-граф-аналитик} (development\_system).
  Исходные данные: граф процесса \texttt{add\_tour} v1.0, модели \texttt{statistic/models.py},
  \texttt{model\_app/tasks.py}.
}

\end{document}
